\documentclass[11pt]{article}
\usepackage[margin=1in]{geometry}
\usepackage{amsmath,amssymb,booktabs}

\newcommand{\E}{\mathbb{E}}
\newcommand{\Var}{\operatorname{Var}}
\newcommand{\Cov}{\operatorname{Cov}}
\newcommand{\sinc}{\operatorname{sinc}}

\title{Kernel error of the ideal encoder:\\finite size, independent noise, shared grid noise}
\author{}
\date{}

\begin{document}
\maketitle

\noindent Three sources make the similarity the memory reads differ from the target Gaussian.
A finite number of waves gives an error that is frozen into the code. Independent noise on each
unit gives an error that is redrawn on every read. Noise in the grid phases is shared by every
wave. Each has a different dependence on the number of waves $M$, the kernel width $r$ and the
arena size $L$. Every check quoted below comes from
\texttt{analysis/hopfield\_probe/kernel\_scaling/kernel\_error\_checks.py} ($L = 1716$, $r = 16$
unless stated).

\section{Setup}

The arena is a torus of side $L$ cells in $d$ dimensions ($d = 2$ for navigation). The code has
$M$ waves. Wave $k$ has an integer frequency $n_k \in \mathbb{Z}^d$, drawn independently from a
Gaussian with standard deviation $\sigma_n = L/(2\pi r)$ per axis and rounded, so its angular
frequency $\omega_k = 2\pi n_k/L$ is approximately $\mathcal{N}(0, I/r^2)$. The code is
\begin{equation}
  z(p) = \frac{1}{\sqrt{M}}\bigl(\cos\omega_k\!\cdot p,\; \sin\omega_k\!\cdot p\bigr)_{k=1}^{M}
  \in \mathbb{R}^{2M}, \qquad |z(p)| = 1 .
\end{equation}
Two positions a displacement $\Delta$ apart have similarity
\begin{equation}
  K_M(\Delta) = z(p)\cdot z(p+\Delta) = \frac{1}{M}\sum_{k=1}^{M}\cos(\omega_k\!\cdot\Delta),
\end{equation}
which does not depend on $p$. The target is
\begin{equation}
  k(\Delta) = \E_\omega\cos(\omega\cdot\Delta) = e^{-|\Delta|^2/2r^2},
\end{equation}
the characteristic function of the Gaussian. The kernel width is set by the spread of the
frequencies alone. $M$ only controls how closely the sum follows its mean.

\section{Error from a finite number of waves}

\subsection{Size at each displacement}

Each term $\cos(\omega_k\cdot\Delta)$ has mean $k(\Delta)$. Its mean square is
$\E\cos^2(\omega\cdot\Delta) = \tfrac12\bigl(1 + k(2\Delta)\bigr) = \tfrac12\bigl(1 + k(\Delta)^4\bigr)$,
so its variance is $\tfrac12(1 - k(\Delta)^2)^2$. The $M$ terms are independent, so
\begin{equation}
  K_M(\Delta) = k(\Delta) + \varepsilon(\Delta), \qquad
  \E\varepsilon = 0, \qquad
  \operatorname{sd}\varepsilon(\Delta) = \frac{1 - e^{-|\Delta|^2/r^2}}{\sqrt{2M}} .
  \label{eq:sdM}
\end{equation}
The error is zero at the goal, where $K_M(0) = 1$ exactly, and reaches $1/\sqrt{2M}$ beyond about
$2r$. The same argument gives the covariance
\begin{equation}
  \Cov\bigl(\varepsilon(\Delta), \varepsilon(\Delta')\bigr)
  = \frac{1}{2M}\bigl[k(\Delta-\Delta') + k(\Delta+\Delta')\bigr] - \frac{1}{M}k(\Delta)k(\Delta') .
  \label{eq:covM}
\end{equation}

\noindent Check (1D, 400 seeds, $M = 512$): at $\Delta = 8, 16, 32, 400$ the measured standard
deviations are $0.0072, 0.0206, 0.0309, 0.0299$ against $0.0069, 0.0198, 0.0307, 0.0312$ from
\eqref{eq:sdM}.

\subsection{Near the goal}

Expanding the cosines, $K_M(\delta) = 1 - \tfrac12\,\delta^\top C\,\delta + O(|\delta|^4)$ with the
curvature matrix
\begin{equation}
  C = \frac{1}{M}\sum_k \omega_k\omega_k^\top, \qquad \E C = I/r^2 .
\end{equation}
Since $\omega_a^2 r^2$ is $\chi^2_1$ with variance $2$, each diagonal entry has relative standard
deviation $\sqrt{2/M}$. So the hill has the right shape on average but its width along any
direction is off by a random fraction $\sqrt{2/M}$ (measured $0.479, 0.250, 0.063$ at $M = 8, 32,
512$ against $0.500, 0.250, 0.0625$). The gradient near the goal is $-C\delta$ instead of
$-\delta/r^2$, wrong by a fraction of order $\sqrt{2/M}$. The error never reverses the slope near
the goal unless $M$ is tiny.

\subsection{Far from the goal}

Beyond about $3r$, $k \approx 0$ and \eqref{eq:covM} reduces to $k(\Delta - \Delta')/2M$ (the second
term only reflects $K_M(\Delta) = K_M(-\Delta)$). For large $M$ the central limit theorem makes
$\varepsilon$ a smooth stationary Gaussian field with variance $\sigma_\varepsilon^2 = 1/2M$, correlation
$k(\tau)$, correlation length $r$, and gradient variance $\sigma_\varepsilon^2/r^2$ per axis.

The largest far value follows from the expected number of excursions of the field above a level
$u\sigma_\varepsilon$. Count both signs and use the evenness of $K_M$, so the far region has volume
$V_1 = L - 6r$ in 1D and $V_2 = L^2 - \pi(3r)^2$ in 2D. The Rice formula (1D) and the Euler
characteristic of the excursion set (2D) give
\begin{equation}
  \mu_1(u) = \frac{V_1}{2\pi r}\,e^{-u^2/2}, \qquad
  \mu_2(u) = \frac{V_2}{(2\pi)^{3/2} r^2}\,u\,e^{-u^2/2} .
\end{equation}
Let $u^*$ solve $\mu(u^*) = 1$. Excursions above $u^*$ are rare and close to Poisson, so the
maximum is Gumbel distributed and
\begin{equation}
  A \equiv \E\max_{|\Delta| > 3r}|K_M(\Delta)| \approx \frac{u^* + \gamma/u^*}{\sqrt{2M}},
  \qquad \gamma = 0.577 .
  \label{eq:alias}
\end{equation}
For $L = 1716$ and $r = 16$, $u^* = 2.36$ in 1D and $3.99$ in 2D.

\begin{center}
\begin{tabular}{lrrr}
  \toprule
  & $M$ & predicted $A$ & measured mean \\
  \midrule
  1D & 32 & 0.326 & 0.321 \\
  1D & 128 & 0.163 & 0.162 \\
  1D & 512 & 0.081 & 0.082 \\
  1D & 2048 & 0.041 & 0.041 \\
  2D & 32 & 0.517 & 0.498 \\
  2D & 512 & 0.129 & 0.128 \\
  \bottomrule
\end{tabular}
\end{center}

\noindent To leading order $u^* \approx \sqrt{2\ln N}$ with $N \sim (L/r)^d$ effectively independent
far values, so
\begin{equation}
  A \approx \sqrt{\frac{d\,\ln(L/r)}{M}}, \qquad
  M \approx \frac{d\,\ln(L/r)}{A^2} \quad\text{for a target alias } A .
\end{equation}
The arena size enters only through a logarithm. The kernel width enters the same way.

With very few waves the Gaussian approximation fails and the code makes copies of the goal. A far
point is a near copy if all $M$ angles fall within $\epsilon$ of zero, which has probability
$(\epsilon/\pi)^M$. Among $N$ far points the closest copy has $\epsilon \approx \pi N^{-1/M}$: about
$0.03$ rad for $M = 2$ (four units) and $0.3$ rad for $M = 4$ in 2D, which means similarity close
to $1$ far from the goal (an estimate, not checked).

\subsection{Edge of the monotone region}

A climber gets past a point as long as the true slope beats the slope of the frozen error.
Along a ray the true slope is $(\rho/r)\,e^{-\rho^2/2}$ with $\rho = |\Delta|/r$. Far from the goal
the gradient of $\varepsilon$ has standard deviation $1/(r\sqrt{2M})$. Setting the two equal gives
the edge of the monotone region,
\begin{equation}
  \rho\,e^{-\rho^2/2} = \frac{1}{\sqrt{2M}}, \qquad \rho^2 \approx \ln(2M) + 2\ln\rho ,
\end{equation}
which is $39, 44, 48, 52$ cells for $M = 32, 128, 512, 2048$ at $r = 16$. This one-standard-deviation
criterion matches the 10th percentile of where a 1D profile first stops decreasing ($38, 43, 47, 51$
cells over 300 seeds; medians $53, 55, 60, 61$). In 2D the nearest trap, a strict local maximum
other than the goal, lies at $41$ to $58$ cells at the closest over 8 seeds (medians $55$ to $69$).
The criterion is scale free, so the edge is proportional to $r$ (not checked at other $r$). It grows
with $M$ only as $\sqrt{\ln M}$: 64 times more waves move the early edge from $38$ to $51$ cells.

\subsection{Comparison with an exact kernel}

The exact kernel has Fourier weights $P_n \propto e^{-r^2|\omega_n|^2/2}$ on all integer
frequencies. Its effective number of dimensions is the participation ratio
\begin{equation}
  D_{\text{exact}} = \frac{(\sum_n P_n)^2}{\sum_n P_n^2}
  = \bigl(4\pi\sigma_n^2\bigr)^{d/2} = \Bigl(\frac{L}{\sqrt{\pi}\,r}\Bigr)^{d},
\end{equation}
which in 2D is $L^2/\pi r^2 \approx 3{,}700$ at $r = 16$. That is the arena's area divided by the
area of one field, the number of place cells of radius $r$ needed to tile it. Random waves do not
reproduce the kernel exactly. They reproduce it up to the alias $A$, and that costs $\ln L$
instead of $L^d$. With unequal weights $w_k$ the far-field results hold with $M$ replaced by
$n_{\text{eff}} = (\sum_k w_k)^2/\sum_k w_k^2$.

\section{Independent noise on each unit}

Each unit carries independent Gaussian noise of standard deviation $s$ relative to its amplitude.
After the $1/\sqrt{M}$ normalisation,
\begin{equation}
  \tilde z(p) = z(p) + \xi, \qquad \xi \sim \mathcal{N}\bigl(0, (s^2/M)\,I_{2M}\bigr).
\end{equation}

\subsection{Error in one similarity read}

With the goal stored clean,
\begin{equation}
  \tilde K = \tilde z(p)\cdot z(q) = K_M(p - q) + \xi\cdot z(q), \qquad
  \xi\cdot z(q) \sim \mathcal{N}(0, s^2/M),
\end{equation}
because $|z(q)| = 1$. The error is $s/\sqrt{M}$ at every displacement. It does not depend on
$r$ or $L$. If the goal is also noisy the variance is $2s^2(1+s^2)/M$. Noise in a stored goal is
drawn once, so it acts as an extra frozen error of size $s/\sqrt{M}$ everywhere. Normalising
$\tilde z$ rescales every read by the same factor and does not move any maximum.

\subsection{Position error of a decoder}

The likelihood of $p$ is $\exp(-M|\tilde z - z(p)|^2/2s^2)$. Since
$\sum_i \partial_a z_i\,\partial_b z_i = C_{ab}$, the Fisher information is
\begin{equation}
  J = \frac{M}{s^2}\,C \approx \frac{M}{s^2 r^2}\,I, \qquad
  \sigma_p \ge \frac{s\,r}{\sqrt{M}} \quad\text{per axis.}
\end{equation}
Precision is set by the curvature of the kernel at its peak. At $r = 16$, $M = 512$, $s = 0.1$
this is $0.07$ cells. A decoder confuses the position with a far alias only when
$s \gtrsim \sqrt{M(1 - A)/(4\ln N)}$, about $3.5$ here, so global errors play no role.

\subsection{Spread of an agent that climbs similarity}

A greedy agent reads $\tilde K$ at $p-1$, $p$, $p+1$, each with fresh noise of standard deviation
$\eta = s/\sqrt{M}$, and moves to the largest. At distance $d$ from the goal the three true
values are $K(d) - g\,i$ for $i = -1, 0, 1$ (outward positive), with slope $g = d/r^2$. For three
Gaussian values, $\partial P(\text{argmax}=i)/\partial\mu_i = 1/2\sqrt{\pi}$ and
$\partial P(\text{argmax}=i)/\partial\mu_j = -1/4\sqrt{\pi}$ for $j\ne i$, so for $g \ll \eta$
\begin{equation}
  \E[\text{step}] = -\beta\,\frac{g}{\eta}, \qquad \beta = \frac{3}{2\sqrt{\pi}} = 0.846,
  \qquad \Var[\text{step}] \approx \tfrac23 .
\end{equation}
The distance is then an autoregressive process $d_{t+1} = (1-\kappa)\,d_t + \text{noise}$ with
$\kappa = \beta/(r^2\eta)$, whose stationary variance is
\begin{equation}
  \langle d^2\rangle = \frac{2/3}{2\kappa - \kappa^2}
  \approx \frac{r^2 s}{3\beta\sqrt{M}}, \qquad
  d_{\text{rms}} \approx 0.63\; r\,\sqrt{s}\;M^{-1/4} .
\end{equation}
The spread depends on $r$, $s$ and $M$ only through $r^2 s/\sqrt{M}$. Check (1D, eight
combinations of $r$, $M$, $s$): at $r^2 s/\sqrt{M} = 2.26, 4.53, 9.05$ the prediction is $1.05,
1.40, 1.93$ cells and the measured spreads are $1.01$, $1.38$ to $1.42$, and $1.89$. Averaging
$T$ reads per step replaces $s$ by $s/\sqrt{T}$.

The climber does worse than the decoder by a factor of about $0.63\,M^{1/4}/\sqrt{s}$. The decoder
uses all $2M$ directions in which the code moves, to first order in the step. The climber reads
one scalar, which is second order in the step at the peak.

\section{Shared noise in the grid phases}

If the code is read out from grid modules, noise there is shared by every wave.

\subsection{Setup}

Work on one axis; in 2D the phase errors add over the axes and the attenuation factors below
multiply. The modules have periods $\lambda_\ell \in \{11, 12, 13\}$ with
$L = \prod_\ell\lambda_\ell = 1716$. Module $\ell$ reports position $p + e_\ell$ (in cells), that
is, phase $\psi_\ell = 2\pi(p + e_\ell)/\lambda_\ell$. Wave $k$ forms
$\tilde\theta_k = \sum_\ell j_{k\ell}\,\psi_\ell$ with integers satisfying
$\sum_\ell j_{k\ell}L/\lambda_\ell \equiv n_k \pmod L$, for example
$j = (6n, -n, 7n) \bmod (11, 12, 13)$. Write $u_{k\ell} = j_{k\ell}/\lambda_\ell$ and
\begin{equation}
  \sum_\ell u_{k\ell} = \frac{n_k}{L} + m_k, \qquad m_k \in \mathbb{Z},
\end{equation}
where $m_k$ depends on which integer represents each residue. At an integer position with no
error, $\theta_k = 2\pi p\sum_\ell u_{k\ell} \equiv \omega_k p \pmod{2\pi}$ for any choice. With
error,
\begin{equation}
  \tilde\theta_k = \omega_k p + \phi_k, \qquad \phi_k = 2\pi\sum_\ell u_{k\ell}\,e_\ell .
\end{equation}

\subsection{Consistent error and the choice of representatives}

If every module is off by the same $x$, then $\phi_k = \omega_k x + 2\pi m_k x$. The first term
moves the hill by $x$, so the code inherits the grid's own position error and nothing more. The
second term gives each wave a different extra phase, which scrambles the code, unless $m_k = 0$.
So the representatives should satisfy
\begin{equation}
  \sum_\ell j_{k\ell}\,\frac{L}{\lambda_\ell} = n_k \quad\text{exactly},
\end{equation}
which is always possible: choose $|j| \le \lambda/2$ for two modules and solve for the third. This
is also the condition for $z$ to be a smooth function of continuous position, since moving by
$x$ between lattice points is exactly a consistent error. With $m_k \ne 0$, wave $k$ oscillates
about $m_k$ times per cell between lattice points while looking slow on the lattice.

\noindent Check ($M = 512$), peak similarity after a consistent error $x$:
\begin{center}
\begin{tabular}{lrrr}
  \toprule
  representatives & $x = 0.10$ & $x = 0.25$ & $x = 0.50$ \\
  \midrule
  $j \in [0, \lambda-1]$ & 0.62 & $-0.35$ & $-0.15$ \\
  $|j| \le \lambda/2$ & 0.93 & 0.64 & 0.28 \\
  $\sum_\ell j_\ell L/\lambda_\ell = n$ & 1.00 & 1.00 & 1.00 \\
  \bottomrule
\end{tabular}
\end{center}

\subsection{Mismatched modules}

With $m_k = 0$, solve for module 13: $u_{k,13} = n_k/L - u_{k,11} - u_{k,12}$. Then
\begin{equation}
  \phi_k = \omega_k e_{13} + 2\pi\bigl(u_{k,11}\,x + u_{k,12}\,y\bigr),
  \qquad x = e_{11} - e_{13}, \quad y = e_{12} - e_{13} .
\end{equation}
Only the differences between modules scramble the code. The similarity to a clean goal is
\begin{equation}
  \tilde K(p) = \frac{1}{M}\sum_k \cos\bigl(\omega_k(p + e_{13}) + 2\pi(u_{k,11}x + u_{k,12}y)\bigr).
\end{equation}
For a broad frequency distribution the residues are close to uniform, so $u_{k,11}$ and $u_{k,12}$
are close to independent and uniform on $[-\tfrac12, \tfrac12]$ (measured: $12\Var u_{11} = 0.91$,
$12\Var u_{12} = 1.04$, correlation $0.11$). Then $\E\cos(2\pi u x) = \sinc x = \sin(\pi x)/\pi x$,
the sine terms average to zero, and
\begin{equation}
  \tilde K(p) \approx a\,k(p + e_{13}), \qquad
  a = \sinc(x)\,\sinc(y) \approx e^{-\pi^2(x^2 + y^2)/6} .
  \label{eq:atten}
\end{equation}
Check (random $e_\ell$ with standard deviation $\sigma_e$, 300 draws): mean peak $0.930, 0.750,
0.616$ against $\sinc\,\sinc$ means $0.932, 0.757, 0.615$ at $\sigma_e = 0.1, 0.2, 0.3$. At
$\sigma_e = 0.5$ the hill is low enough that the maximum picks up the far field ($0.39$ against
$0.35$).

Properties of \eqref{eq:atten}:
\begin{itemize}
  \item Neither $M$ nor $r$ appears. The factor $a$ multiplies the whole hill. The only role of
  $r$ is to make the residues uniform, which needs $\sigma_n \gtrsim \lambda$, that is
  $r \lesssim L/2\pi\lambda \approx 21$ here.
  \item It is a bias in the height of the hill, not a fluctuation, so averaging reads does not
  remove it while the modules disagree.
  \item $a = 0$ at a mismatch of one cell. That mismatch is a valid code for another position, so
  the read sees only far-field similarity.
  \item The goal stops being the highest point when $a < A$, that is when
  $x^2 + y^2 > (6/\pi^2)\ln(1/A)$, summed over axes in 2D. For $M = 512$ ($A = 0.13$) this is
  about $1.25$ cells$^2$. More waves help only through $\ln(1/A)$, which grows as
  $\tfrac12\ln M$.
\end{itemize}

For any representatives, the variance of $\phi$ over waves is $4\pi^2 e^\top U e$ with
$U = \Cov_k(u)$. With $m_k = 0$, $U$ has the all-ones vector as an approximate null vector, so
consistent errors drop out. Symmetric representatives have a smaller $U$ overall ($12\operatorname{tr}U
= 2.86$ against $4.12$) and do a little better when the modules err independently (mean peak
$0.95, 0.84, 0.69$ against $0.93, 0.75, 0.62$ at $\sigma_e = 0.1, 0.2, 0.3$), but they fail
under consistent errors. Coupled modules mostly drift together (Mosheiff and Burak, 2019) and
continuous movement is itself a consistent change, so $m_k = 0$ is the right choice.

\section{Summary}

\begin{center}
\small
\begin{tabular}{p{2.3cm}p{4.2cm}p{1.6cm}p{2.6cm}p{2.4cm}}
  \toprule
  source & error in $K$ & $M$ & $r$, $L$ & where, when \\
  \midrule
  finite $M$ &
  sd $(1 - k^2)/\sqrt{2M}$; largest far value $(u^* + \gamma/u^*)/\sqrt{2M}$ &
  $1/\sqrt{M}$ &
  through $\ln(L/r)$ only &
  zero at the goal, full beyond $2r$; frozen \\[4pt]
  independent noise &
  $s/\sqrt{M}$ per read &
  $1/\sqrt{M}$ &
  none; position error $s r/\sqrt{M}$ (decoder), $0.63\,r\sqrt{s}M^{-1/4}$ (climber) &
  same everywhere; fresh each read \\[4pt]
  shared grid noise &
  factor $\sinc(x)\sinc(y)$ on the whole hill &
  none, except through $A$ &
  none while residues are uniform &
  whole hill; lasts while modules disagree \\
  \bottomrule
\end{tabular}
\end{center}

\end{document}
